% Strategic Management - Week 10 Master Vocabulary Bank
% Compile: pdflatex sm-master-vocabulary-word-bank.tex

\documentclass[10pt,a4paper]{article}
\usepackage[utf8]{inputenc}
\usepackage[T1]{fontenc}
\usepackage[margin=1.1cm]{geometry}
\usepackage{xcolor}
\usepackage{fancyhdr}
\usepackage{titlesec}
\usepackage{enumitem}
\usepackage{tabularx}
\usepackage{lmodern}

\definecolor{pinkmain}{RGB}{205,38,124}
\definecolor{pinksoft}{RGB}{255,239,248}
\definecolor{pinkdark}{RGB}{126,20,77}
\definecolor{textdark}{RGB}{38,38,38}

\setlength{\parindent}{0pt}
\setlength{\parskip}{0.35em}
\setlist[itemize]{leftmargin=1.2em,itemsep=2pt,topsep=2pt}

\titleformat{\section}{\normalsize\bfseries\color{pinkdark}}{}{0pt}{}
\titleformat{\subsection}{\small\bfseries\color{pinkdark}}{}{0pt}{}

\pagestyle{fancy}
\fancyhf{}
\fancyfoot[L]{\scriptsize Strategic Management | Week 10 | Master Vocabulary}
\fancyfoot[R]{\scriptsize \thepage}

\newcommand{\vocab}[2]{\item \textbf{\textcolor{pinkmain}{#1}} -- #2}

\begin{document}
\begin{center}
{\Large\bfseries\color{pinkdark} Strategic Management Master Vocabulary Bank}\\
{\small\color{textdark} Week 10 consolidation sheet (Units 0-6 + case language)}
\end{center}

\colorbox{pinksoft}{\parbox{\dimexpr\linewidth-2\fboxsep}{\small
\textbf{Source scope:} terms consolidated from the `strategic-management/` folder materials (unit summaries, class notes, and study guides through Week 9).}}

\section*{Unit 0-1 Foundations}
\begin{itemize}
\vocab{Strategic management}{Field focused on long-term success through analysis, decisions, and implementation.}
\vocab{Strategy}{Long-term direction plus coordinated actions and resource allocation.}
\vocab{Strategic decision}{High-impact, hard-to-reverse choice under uncertainty and complexity.}
\vocab{Sustained success}{Performance advantage maintained over time, not one-time luck.}
\vocab{Corporate strategy}{Business scope choices: where to compete and portfolio logic.}
\vocab{Competitive strategy}{How to compete inside a business versus rivals.}
\vocab{Functional strategy}{How functions (HR, marketing, operations, finance, R\&D) support higher-level strategy.}
\vocab{Strategic target}{Industry-wide scope or segment focus of competition.}
\vocab{Deliberate strategy}{Planned strategy process: analysis, formulation, implementation, control.}
\vocab{Emergent strategy}{Strategy pattern that arises from adaptation and changing conditions.}
\vocab{Integrative strategy process}{View combining intended plans with emergent adjustments.}
\vocab{Limited rationality}{Decision constraint from incomplete information and cognitive limits.}
\vocab{Top Management Team (TMT)}{Senior executives coordinating cross-functional strategic execution.}
\vocab{Board of Directors}{Governance body for advice, supervision, and owner-manager control bridge.}
\vocab{Scarce resources}{Limited assets (customers, talent, capital) firms compete to secure.}
\vocab{Intended vs realized strategy}{Difference between planned strategic direction and actual outcomes.}
\end{itemize}

\section*{Unit 2 Objectives, Values, Governance}
\begin{itemize}
\vocab{Mission}{Present identity and purpose of the firm: what business it is in.}
\vocab{Vision}{Long-term aspiration of what the firm aims to become.}
\vocab{Strategic objective}{Specific target with attribute, metric, level, and timeframe.}
\vocab{Yardstick}{Measurement indicator used to track objective attainment (e.g., ROE).}
\vocab{Firm performance}{Degree to which economic, market, and social goals are achieved.}
\vocab{ROA}{Return on assets; efficiency of profit generation from total assets.}
\vocab{ROE}{Return on equity; profitability from shareholders' invested capital.}
\vocab{Shareholder value}{Market value created for owners (price appreciation plus distributions).}
\vocab{Opportunity cost}{Return foregone by not choosing the best alternative investment.}
\vocab{Cost of equity}{Minimum return shareholders require given risk profile.}
\vocab{Stakeholder}{Actor affected by or influencing firm decisions (employees, customers, state, etc.).}
\vocab{Shareholder primacy}{View that firm objective is maximizing owner wealth.}
\vocab{Stakeholder theory}{View that firms should balance value across multiple stakeholder groups.}
\vocab{Triple bottom line}{Economic, social, and environmental performance frame.}
\vocab{Corporate governance}{Mechanisms aligning managerial behavior with ownership objectives.}
\vocab{Agency problem}{Conflict where managers' interests diverge from owners' interests.}
\vocab{Incentive alignment}{Compensation design linking manager rewards to owner value creation.}
\vocab{Corporate Social Responsibility (CSR)}{Voluntary actions addressing social/environmental stakeholder demands.}
\vocab{Business ethics}{Moral principles guiding strategic decisions and business conduct.}
\vocab{Corporate values}{Core beliefs and standards shaping behavior, culture, and decisions.}
\end{itemize}

\section*{Unit 3 External Analysis}
\begin{itemize}
\vocab{Business environment}{External conditions that influence firm behavior and results but cannot be controlled directly.}
\vocab{Macro-environment}{Broad contextual forces captured by PESTEL analysis.}
\vocab{Industry environment}{Competitive arena where firms, buyers, suppliers, entrants, and substitutes interact.}
\vocab{PESTEL}{Political, Economic, Social, Technological, Ecological, Legal framework for macro analysis.}
\vocab{Industry attractiveness}{Expected ability of an industry to sustain above-normal returns.}
\vocab{Five Forces}{Porter model of rivalry, entry, substitutes, buyer power, supplier power.}
\vocab{Rivalry among existing firms}{Intensity of competition among current players in an industry.}
\vocab{Threat of new entrants}{Likelihood that new competitors can enter and erode profits.}
\vocab{Barriers to entry}{Obstacles that make market entry difficult or costly.}
\vocab{Natural barriers}{Structural barriers like scale economies, know-how, capital intensity.}
\vocab{Artificial barriers}{Institutional barriers like regulation, tariffs, legal requirements.}
\vocab{Threat of substitutes}{Risk from alternative offerings satisfying the same customer need.}
\vocab{Disruptive technology}{Technology that changes consumption patterns and industry economics.}
\vocab{Buyer bargaining power}{Customer ability to force lower prices or better terms.}
\vocab{Supplier bargaining power}{Input provider ability to raise prices or constrain quality/supply.}
\vocab{Forward integration threat}{Risk suppliers enter downstream and compete directly.}
\vocab{Backward integration threat}{Risk buyers internalize upstream activities.}
\vocab{Industry effect}{Performance portion explained by structural industry conditions.}
\vocab{Firm effect}{Performance portion explained by firm-specific resources and strategy.}
\vocab{Cluster}{Geographic concentration of related firms, suppliers, and institutions.}
\vocab{Strategic group}{Subset of firms in an industry with similar strategic positioning.}
\end{itemize}

\section*{Unit 4 Internal Analysis (RBV, Value Chain, SWOT)}
\begin{itemize}
\vocab{Internal diagnosis}{Assessment of strengths and weaknesses from inside the firm.}
\vocab{Firm identity}{Profile of size, age, scope, ownership, and activity type.}
\vocab{Value chain}{Decomposition of activities through which value is created and delivered.}
\vocab{Primary activities}{Inbound logistics, operations, outbound logistics, marketing/sales, service.}
\vocab{Supporting activities}{Infrastructure, HRM, technology development, procurement.}
\vocab{Value system}{Linked value chains of firm, suppliers, channels, and customers.}
\vocab{Horizontal linkages}{Coordination and optimization links across internal activities.}
\vocab{Vertical linkages}{Coordination links with suppliers or buyers in adjacent stages.}
\vocab{Resource-Based View (RBV)}{Theory that heterogenous resources and capabilities explain performance differences.}
\vocab{Resource}{Asset the firm controls (tangible, intangible, or human).}
\vocab{Capability}{Coordinated ability to deploy resources through routines and processes.}
\vocab{Core competence}{Capability performed better than rivals and central to advantage.}
\vocab{Organizational routine}{Regular, repeatable pattern of coordinated organizational action.}
\vocab{Tangible resource}{Physical/financial asset usually measurable in accounts.}
\vocab{Intangible resource}{Non-physical asset (brand, culture, reputation, technology knowledge).}
\vocab{Human resource}{People-based knowledge, motivation, and know-how.}
\vocab{Scarcity (RBV)}{Resource not broadly available among competitors.}
\vocab{Relevance (RBV)}{Resource fit with key success factors of the industry.}
\vocab{Durability}{Ability of a resource/capability to retain value over time.}
\vocab{Inimitability}{Difficulty for competitors to copy the resource/capability.}
\vocab{Tradeability}{Ease with which a resource can be bought/sold across firms.}
\vocab{Substitutability}{Extent to which alternatives can replace a strategic resource.}
\vocab{Complementarity}{Extra value generated by combining resources/capabilities together.}
\vocab{Rent appropriation}{Who captures economic value generated by strategic resources.}
\vocab{VRIN}{Valuable, Rare, Inimitable, Non-substitutable criteria for sustained advantage.}
\vocab{VRIO}{VRIN plus organization to capture value from resources/capabilities.}
\vocab{SWOT}{Synthesis matrix of strengths, weaknesses, opportunities, threats.}
\vocab{Strengths}{Internal attributes supporting competitive success.}
\vocab{Weaknesses}{Internal limitations harming performance potential.}
\vocab{Opportunities}{External conditions that can be exploited for strategic gain.}
\vocab{Threats}{External conditions that can damage performance or position.}
\end{itemize}

\section*{Unit 5 Competitive Strategy}
\begin{itemize}
\vocab{Competitive parity}{No meaningful relative advantage versus industry average rivals.}
\vocab{Customer value}{Perceived benefit received by buyers from the firm's offer.}
\vocab{Margin}{Difference between price and cost captured by the firm.}
\vocab{Cost leadership}{Lower-cost position for comparable offering quality.}
\vocab{Cost driver}{Structural source of lower costs (scale, learning, process, design, inputs).}
\vocab{Economies of scale}{Lower average cost as output increases.}
\vocab{Learning economies}{Efficiency gains from repetition, experience, and routine improvement.}
\vocab{X-inefficiency}{Performance losses from weak discipline/motivation inside organizations.}
\vocab{Differentiation}{Perceived uniqueness valued by customers and potentially premium-priced.}
\vocab{Price premium}{Higher price justified by distinctive perceived value.}
\vocab{Differentiation driver}{Source of uniqueness (features, service, reputation, legitimacy, timing).}
\vocab{Adequacy condition}{Contextual condition required for strategy fit and effectiveness.}
\vocab{Commodity market}{Low-differentiation market where price competition dominates.}
\vocab{Stuck in the middle}{Strategic inconsistency with no strong cost edge or differentiation edge.}
\vocab{Strategy clock}{Eight-position model mapping perceived value against price.}
\vocab{Hybrid position}{Clock position combining moderate differentiation with competitive pricing.}
\vocab{Focused differentiation}{Premium strategy for a narrow target segment.}
\vocab{Industry life cycle}{Stage model: emergence, growth, maturity, decline.}
\vocab{Shakeout}{Stage where weaker firms exit and competition concentrates.}
\vocab{Harvest strategy}{Cash-maximizing approach in declining businesses/markets.}
\vocab{Rapid withdrawal}{Fast exit from unattractive/declining positions to preserve value.}
\end{itemize}

\section*{Unit 6 Corporate Strategy Development}
\begin{itemize}
\vocab{Where to compete}{Core corporate strategy question about scope and business domain.}
\vocab{Firm scope}{Set of businesses/activities in which the company competes.}
\vocab{Corporate development}{Scope changes through specialization, diversification, integration, restructuring.}
\vocab{Firm growth}{Increase/decrease in size metrics (assets, outputs, sales, profits).}
\vocab{Internal growth}{Expansion using internal resources and organic capability building.}
\vocab{Mergers and acquisitions (M\&A)}{External growth via ownership combinations/purchases.}
\vocab{Strategic alliance}{Cooperative arrangement between firms without full ownership integration.}
\vocab{Specialization}{Concentration on core business domains and known capabilities.}
\vocab{Diversification}{Expansion into additional businesses (related or unrelated).}
\vocab{Related diversification}{Expansion into businesses sharing capabilities or value chain links.}
\vocab{Unrelated diversification}{Expansion into businesses with limited operational similarity.}
\vocab{Economies of scope}{Cost/value gains from sharing resources across businesses.}
\vocab{Vertical integration}{Internalization of multiple sequential value-chain stages.}
\vocab{Vertical disintegration}{Outsourcing/splitting value-chain stages instead of internalizing them.}
\vocab{Opportunism risk}{Risk of self-interested behavior by external transacting parties.}
\vocab{Restructuring}{Redefinition of scope including refocus, turnaround, or divestment.}
\vocab{Divestment}{Disposal or withdrawal from a business unit/activity.}
\vocab{Portfolio restructuring}{Scope redesign via both business exits and entries.}
\vocab{Re-float (turnaround)}{Actions to recover and stabilize underperforming businesses.}
\vocab{Winding-up}{Full closure of operations and sale of remaining assets.}
\end{itemize}

\section*{Case and Exam Language Bank}
\begin{itemize}
\vocab{Repeatable business model execution}{Consistent system-level delivery across units/locations.}
\vocab{Third place}{Positioning concept: social space between home and work.}
\vocab{Advantage renewal}{Continuous adaptation required to sustain, not just defend, advantage.}
\vocab{Transferability}{Extent to which strategy success in one context can be replicated in another.}
\vocab{Turnaround}{Coordinated managerial actions to reverse decline and restore performance.}
\vocab{System coherence}{Alignment between positioning, capabilities, operations, and scope choices.}
\vocab{Case diagnosis}{Evidence-led identification of strategic issue, cause, and options.}
\vocab{Framework mapping}{Using course models to interpret case facts and justify recommendations.}
\vocab{Strategic implication}{Action consequence derived from analysis findings.}
\vocab{Boundary condition}{Condition under which a strategic claim holds or fails.}
\end{itemize}

\section*{Fast Oral Template (Vocabulary in Use)}
\colorbox{pinksoft}{\parbox{\dimexpr\linewidth-2\fboxsep}{\small
\textbf{Template:} ``Using \textbf{[framework]}, the key issue is \textbf{[problem]}. The firm has \textbf{[resource/capability or market condition]}, which creates \textbf{[advantage/risk]}. Therefore, it should \textbf{[action]}, because this improves \textbf{[fit, sustainability, or scope coherence]}.''}}

\end{document}
